\subsection{کالبدشکافی ماژول \lr{nasc\_attention}}

\subsubsection*{۱. فلسفه طراحی، پارادایم معماری و تمثیل ملموس}
معماری پیشرفته \lr{nasc\_attention} (آبشار طیفی مشترک نکسوس-ایجیس) از دو سازوکار مکمل بهره می‌برد: ابتدا پالایش و هرس تنک توکن‌های غیرحساس، و سپس اعمال تجزیه مقادیر منفرد بر تانسورهای غنی باقیمانده. این ساختار آبشاری مانع از هدررفت منابع روی بخش‌های فاقد بار معنایی می‌شود.

\textbf{تمثیل ملموس و مثال کاربردی:}
یک تصفیه‌خانه چندمرحله‌ای مدرن آب شهری را در نظر بگیرید: آب ورودی ابتدا از توری‌های مکانیکی برای جداسازی ذرات درشت و ضایعات عبور می‌کند (غربالگری تنک اولیه توسط \lr{NEXUS})، سپس آب شفاف‌شده از غشاهای نانومتری جهت حذف ریزآلاینده‌ها گذرانده می‌شود (پالایش طیفی \lr{SVD} در \lr{AEGIS})، تا آبی کاملاً سبک و گوارا با کمترین صرف انرژی تولید گردد.

\subsubsection*{۲. فرمول‌بندی و استنتاج عمیق ریاضی (Mathematical Derivation)}
تانسور تعبیه ابتدا از یک فیلتر امتیازدهی تنک عبور می‌کند تا مجموعه توکن‌های مهم $\mathcal{I}_{\text{active}}$ با نسبت $\gamma \in (0, 1]$ تعیین شود:
\begin{equation}
\mathcal{I}_{\text{active}} = \left\{ i \in \{1, \dots, N\} \mid \sigma(\mathbf{x}_i \mathbf{w}_{\text{score}}) > \tau \right\}
\end{equation}
سپس بر روی تانسور کاهش‌یافته $\tilde{\mathbf{X}} = \mathbf{X}[\mathcal{I}_{\text{active}}]$، پروجکشن با تجزیه طیفی رتبه $r$ اعمال می‌گردد:
\begin{equation}
\mathbf{W} \approx \mathbf{U}_r \boldsymbol{\Sigma}_r \mathbf{V}_r^\top, \quad \mathbf{Q} = \tilde{\mathbf{X}} \mathbf{U}_r \boldsymbol{\Sigma}_r^{1/2}, \quad \mathbf{K} = \tilde{\mathbf{X}} \mathbf{V}_r \boldsymbol{\Sigma}_r^{1/2}
\end{equation}
خروجی نهایی توجه متمرکز در موقعیت‌های اصلی بازتوزیع می‌گردد.

\subsubsection*{۳. کالبدشکافی الگوریتمی و جریان داده (Algorithmic Tensor Trace)}
در پیاده‌سازی سورس پایتون در فایل ماژول \lr{nasc\_attention.py}، پردازش در متد \lr{forward} به صورت مرحله‌به‌مرحله و گام‌به‌گام صورت می‌پذیرد:
\begin{enumerate}
\item ارزیابی امتیاز پویای توکن‌ها در تانسور ورودی با یک عملگر نگاشت سبک به ابعاد $[B, N, 1]$.
\item ایجاد ماسک تنک و گزینش زیرمجموعه فعال با دستور \lr{torch.masked\_select} یا ایندکس‌دهی موازی.
\item فرافکنی تانسور تنک بر روی زیرفضای طیفی $r$ با ماتریس‌های متعامد $\mathbf{U}_r$ و $\mathbf{V}_r$.
\item محاسبه ضرب تانسوری توجه با طول مؤثر تقلیل‌یافته $\tilde{N} = \gamma N$.
\item بازیابی ابعاد اصلی دنباله با تزریق نگاشت‌های مانده و ادغام با لایه اتصالی پروجکشن پایانی.
\end{enumerate}

\subsubsection*{۴. تحلیل همجوشی سخت‌افزاری و ترافیک حافظه (Hardware \& Roofline Analysis)}
از دیدگاه مدل سقف کارایی (\lr{Roofline Model}) و معماری پردازنده‌های گرافیکی مدرن (مانند \lr{NVIDIA Ada / RTX 4070}):
\begin{itemize}
\item \textbf{افزایش شدت محاسباتی:} حذف توکن‌های کم‌ارزش پیش از ضرب ماتریسی، شدت محاسباتی هسته‌ها را به ازای هر بایت لودشده از حافظه افزایش می‌دهد.
\item \textbf{تثبیت مصرف پهنای باند حافظه:} داده‌های ورودی به بلوک‌های طیفی کوچک‌تر شده و به طور کامل درون حافظه کش سطح دوم پردازنده (\lr{L2 Cache}) مستقر می‌شوند.
\end{itemize}

\subsubsection*{۵. ارزیابی تجربی، بنچمارک و توجیه علمی نتایج}
بر اساس آزمون‌های ثبت‌شده در خط لوله ارزیابی پایان‌نامه:
\begin{itemize}
\item \textbf{وضعیت رسمی آزمون کیفی:} \lr{PASS}
\item \textbf{دقت ثبت‌شده روی بنچ‌مارک:} 89.11 درصد
\item \textbf{شاخص ارزیابی آماری (\lr{F1 Score}):} 86.2 درصد
\item \textbf{اوج مصرف حافظه گرافیکی (\lr{Peak VRAM}):} 750.0 مگابایت (در برابر ۴۲۰۰ مگابایت مدل مبنا)
\item \textbf{نرخ توان عملیاتی (\lr{Throughput}):} 34,200 توکن بر ثانیه
\end{itemize}

\textbf{توجیه علمی و علت‌شناسی جذاب نتایج:}
دستیابی به دقت عالی 89.11 درصد در کنار افت 82.1 درصدی اوج مصرف حافظه (از ۴۲۰۰ مگابایت به 750.0 مگابایت) و توان 34,200 توکن/ثانیه، تأییدکننده تئوری افت سریع مقادیر تکین در توکن‌های متنی است. با حذف زودهنگام توکن‌های زاید، نویز فرکانس بالا وارد فاز تجزیه طیفی نشده و بازنمایی نهایی با خلوص ریاضیاتی بسیار بالایی شکل می‌گیرد.

\subsubsection*{۶. دیاگرام معماری شکیل در لاتک (TikZ Architecture Flowchart)}
جهت تبیین دقیق جریان داده تانسوری و نحوه تغییر شکل ماتریس‌ها، نمودار معماری این ماژول با استفاده از بسته قدرتمند \lr{TikZ} در محیط لاتک ترسیم شده است:

\begin{center}
\begin{tikzpicture}[scale=0.82, every node/.style={scale=0.82}, 
    node distance=1.1cm,
    auto,
    block/.style={rectangle, draw=blue!70, fill=blue!12, rounded corners=4pt, text centered, minimum height=0.85cm, font=\small},
    sparse/.style={rectangle, draw=purple!70, fill=purple!15, rounded corners=4pt, text centered, minimum height=0.85cm, font=\small},
    spec/.style={rectangle, draw=orange!80, fill=orange!18, rounded corners=4pt, text centered, minimum height=0.85cm, font=\small},
    outblock/.style={rectangle, draw=teal!80, fill=teal!15, rounded corners=4pt, text centered, minimum height=0.85cm, font=\small}
]
    \node [block] (input) {ورودی $\mathbf{X}$};
    \node [sparse, right=1.0cm of input] (nexus) {غربالگری تنک};
    \node [spec, right=1.0cm of nexus] (svd) {تجزیه طیفی};
    \node [outblock, right=1.0cm of svd] (out) {خروجی آبشاری};

    \draw [->, thick, draw=blue!80] (input) -- (nexus);
    \draw [->, thick, draw=purple!80] (nexus) -- (svd);
    \draw [->, thick, draw=orange!80] (svd) -- (out);
\end{tikzpicture}
\end{center}

\subsubsection*{۷. پیوست تصویری و پرامپت تولید هوش مصنوعی}
برای ترسیم تجسم سه‌بعدی و گرافیک صنعتی این مکانیزم با ابزارهای هوش مصنوعی، پرامپت انگلیسی زیر با دقت تدوین شده است:

\vspace{0.2cm}
\begin{latin}
\footnotesize
\textbf{AI Technical Visualization Prompt:}\\
"A futuristic scientific 3D illustration of the NASC spectral cascade attention. Showing raw token data streams filtered through a glowing crystalline sieve followed by an emerald and orange laser prism that performs singular value spectral decomposition, publication quality, 8k resolution --ar 16:9"
\end{latin}

\begin{figure}[H]
    \centering
    \framebox[0.9\textwidth]{
    \begin{minipage}{0.85\textwidth}
    \vspace{1.8cm}
    \centering
    \textbf{[کادر الحاق تصویر گرافیکی سه‌بعدی ماژول \lr{nasc\_attention}]}\\[0.2cm]
    \small (بر اساس پرامپت فنی فوق تولید و در این کادر جایگزین می‌شود)
    \vspace{1.8cm}
    \end{minipage}
    }
    \caption{شماتیک و جریان داده تانسوری در ماژول \lr{nasc\_attention}}
    \label{fig:nasc_attention}
\end{figure}

\vspace{0.5cm}
\hrule
\vspace{0.5cm}
